%%%%%%%%%%%%%%%%%%%%%%%%%%%%%%%%%%%%%%%%%%%%%%%%%%%%%%%%%%%%%%%%%%%%%%%%%%
%%%%%% Optimization & Linear Algebra for Deep Learning
%%%%%% NPTEL Course Notes (Lectures 5.1--6.6)
%%%%%% Revised May 14 2026 — spacing fixes + handwritten-note observations
%%%%%%%%%%%%%%%%%%%%%%%%%%%%%%%%%%%%%%%%%%%%%%%%%%%%%%%%%%%%%%%%%%%%%%%%%%
\documentclass{beamer}
\hypersetup{pdfpagemode=FullScreen}
\setbeamertemplate{navigation symbols}{}
\usepackage[english]{babel}
\usepackage[latin1]{inputenc}
\usepackage{amsmath,amssymb}
\usepackage[T1]{fontenc}
\usepackage{times}
\usepackage{mathptmx}
\usepackage{graphicx}
\usepackage{booktabs}
\usepackage{hyperref}
\usepackage{xcolor}
\usepackage{mathtools}
\usepackage{algorithm,algorithmic}
\renewcommand{\algorithmicrequire}{\textbf{Input:}}
\renewcommand{\algorithmicensure}{\textbf{Output:}}

\mode<presentation>{
  \usetheme{Madrid}
  \usecolortheme{whale}
}

% Give frames more breathing room: reduce block title padding and use small font
\setbeamertemplate{blocks}[rounded][shadow=false]
\AtBeginEnvironment{block}{\small}
\AtBeginEnvironment{alertblock}{\small}
\AtBeginEnvironment{exampleblock}{\small}
% Shrink all frames slightly so content fits
\setbeamersize{text margin left=6pt, text margin right=6pt}

\setbeamercovered{transparent}
\setbeamertemplate{bibliography item}[text]
\setbeamertemplate{theorems}[numbered]
\setbeamerfont{title}{size=\large}
\setbeamerfont{date}{size=\tiny}
\setbeamertemplate{itemize items}[ball]
\setbeamertemplate{itemize subitem}[circle]
\setbeamertemplate{itemize subsubitem}[triangle]
\setbeamertemplate{caption}[numbered]

%------------------------------------------------------------------
%  TITLE PAGE
%------------------------------------------------------------------
\title[Optimization \& Linear Algebra]{%
  Optimization Algorithms \& Linear Algebra\\[4pt]
  for Deep Learning}

\author[Vatsav Reddy]{%
  \texorpdfstring{%
    {\Large Vatsav Reddy}\\[6pt]
    {\small Under the Guidance of}\\[2pt]
    {\small Prof.\ Praveen Parchuri}%
  }{Vatsav Reddy}}

\institute[IIITH]{%
  {\normalsize International Institute of Information Technology, Hyderabad}}

\date{\today}

\begin{document}

\begin{frame}
  \titlepage
\end{frame}

\begin{frame}[shrink=10]{Contents}
  \tableofcontents
\end{frame}

\AtBeginSection[]{
  \begin{frame}<beamer>
    \frametitle{Current Section}
    \tableofcontents[currentsection]
  \end{frame}
}

%==================================================================
\section{Lec 5.1 \& 5.2 Recap: Gradient Descent \& Contour Maps}
%==================================================================

\begin{frame}[shrink=10]{Gradient Descent Recap}
  \begin{block}{Setup}
    Prediction $\hat{y} = f(x) = \tfrac{1}{1+e^{-(w\cdot x+b)}}$,\quad
    parameter vector $w = \{w^1, w^2, \ldots\}$,\quad
    labels $q = \{q_1,\ldots,q_h\}$,\quad
    $\mathcal{L}(p\|q) = \mathrm{CE}(p,q)$.
  \end{block}

  \medskip

  \begin{block}{Update Rule: $w \leftarrow w - \eta\,\nabla w$}
    \begin{itemize}\setlength\itemsep{3pt}
      \item \textbf{Steep:}  $\frac{\Delta y_1}{\Delta x_1}$ large
            $\Rightarrow$ updates \textbf{large} $\Rightarrow$ \textbf{Fast}
      \item \textbf{Gentle:} $\frac{\Delta y_2}{\Delta x_2}$ small
            $\Rightarrow$ updates \textbf{small} $\Rightarrow$ \textbf{Slow}
    \end{itemize}
  \end{block}
\end{frame}

%------------------------------------------------------------------
\begin{frame}[shrink=10]{Contour Maps (Lec 5.3)}
  \begin{block}{Key Observation}
    The error surface can be visualised in 3-D (over $w$ and $b$ axes)
    or as 2-D contours.
    Moving $3\text{D}\to 2\text{D}\to$ ``combination'' view aids intuition.
  \end{block}

  \begin{columns}[T]
    \column{0.52\textwidth}
    \begin{block}{Reading Contours}
      \begin{itemize}\setlength\itemsep{4pt}
        \item Ellipse/circle boundary $\Rightarrow$ \textbf{same error}
        \item \textbf{Small} gap $\Rightarrow$ steep slope
        \item \textbf{Large} gap $\Rightarrow$ gentle slope
        \item Plateau $\Rightarrow$ very slow movement
      \end{itemize}
    \end{block}

    \column{0.44\textwidth}
    \begin{alertblock}{Key Rule (proportional to gradient)}
      \begin{itemize}\setlength\itemsep{4pt}
        \item Gentle $\to$ smaller gradient
              $\to$ \textbf{slower moment}
        \item Steep $\to$ higher gradient
              $\to$ \textbf{higher moment}
      \end{itemize}
    \end{alertblock}
  \end{columns}
\end{frame}

%==================================================================
\section{Lec 5.4: Momentum-Based Descent}
%==================================================================

\begin{frame}[shrink=10]{Problem with Vanilla Gradient Descent}
  \begin{alertblock}{Problem}
    Vanilla GD takes a lot of time to navigate regions with a
    \textbf{gentle slope} because the gradient is very small there.
  \end{alertblock}

  \bigskip

  \begin{block}{Momentum GD --- Intuition}
    \textit{``If I am repeatedly being asked to move in the same direction,
    then I should probably gain some confidence and start taking
    \textbf{bigger steps} in that direction.''}
  \end{block}

  \bigskip

  \begin{block}{Regularisation Note}
    \begin{itemize}\setlength\itemsep{6pt}
      \item \textbf{L1 (Lasso):}
            $\mathcal{L}_1 = \displaystyle\sum_{i}(y_i-\hat{y}_i)^2
            + \lambda\sum_j|w_j|$ \quad (encourages \emph{sparsity})
      \item \textbf{L2 (Ridge):}
            $\mathcal{L}_2 = \displaystyle\sum_{i}(y_i-\hat{y}_i)^2
            + \lambda\sum_j w_j^2$ \quad (\emph{non-sparse} / shrinkage)
    \end{itemize}
  \end{block}
\end{frame}

%------------------------------------------------------------------
\begin{frame}[shrink=10]{Momentum Update Rule}
  \begin{block}{Update Equations}
    \begin{align*}
      \text{update}_t &= \beta\cdot\text{update}_{t-1}
                        + \eta\,\nabla w_t
                        \qquad\leftarrow\text{current gradient}\\[4pt]
      w_{t+1}        &= w_t - \text{update}_t
                       = w_t - \beta\,\text{update}_{t-1}
                         - \eta\,\nabla w_t
    \end{align*}
    In addition to the current update we also look at the
    \textbf{history of updates}.
  \end{block}

  \medskip

  \begin{exampleblock}{Exponentially Weighted Average --- Derivation ($\text{update}_0=0$)}
    \begin{align*}
      \text{update}_1 &= \eta\nabla w_1\\
      \text{update}_2 &= \beta\eta\nabla w_1 + \eta\nabla w_2\\
      \text{update}_3 &= \beta^2\eta\nabla w_1
                         + \beta\eta\nabla w_2 + \eta\nabla w_3\\
      \vdots\;&
    \end{align*}
    \[
      \boxed{\text{update}_t
        = \beta^{t-1}\eta\nabla w_1 + \cdots + \eta\nabla w_t}
    \]
  \end{exampleblock}
\end{frame}

%------------------------------------------------------------------
\begin{frame}[shrink=10]{Why Momentum Works}
  \begin{block}{Key Insight}
    Even in regions with \textbf{gentle slopes}, momentum-based GD can
    take \textbf{large steps} because the accumulated momentum carries
    it along.
  \end{block}

  \bigskip

  \begin{alertblock}{Is Moving Fast Always Good?}
    No!\ Momentum can take many \textbf{U-turns} before converging.
    Despite these U-turns it still converges \emph{faster} than vanilla GD.
  \end{alertblock}

  \bigskip

  \begin{exampleblock}{Observation from Notes}
    The U-turns arise because momentum overshoots; the gradient
    corrects direction repeatedly.
    Convergence is still guaranteed and faster than plain GD.
  \end{exampleblock}
\end{frame}

%==================================================================
\section{Lec 5.5: Nesterov Accelerated Gradient Descent}
%==================================================================

\begin{frame}[shrink=10]{NAG: Look Before You Leap}
  \begin{columns}[T]
    \column{0.48\textwidth}
    \begin{block}{Momentum (standing here)}
      \begin{enumerate}\setlength\itemsep{4pt}
        \item Stand here
        \item Check slope \emph{here}
        \item Move
      \end{enumerate}
      \smallskip
      \textit{A single data point --- current position only}
    \end{block}

    \column{0.48\textwidth}
    \begin{block}{Nesterov (standing ahead)}
      \begin{enumerate}\setlength\itemsep{4pt}
        \item Peek ahead
        \item Check \textbf{future} slope
        \item Adjust \emph{before} moving
      \end{enumerate}
    \end{block}
  \end{columns}

  \medskip

  \begin{block}{NAG --- Key Difference \& Update Rule}
    Compute the gradient at the \textbf{look-ahead} value of $w$:
    \[
      w_{\text{look-ahead}} = w_t - \beta\,\text{update}_{t-1}
    \]
    (Move according to history first --- you might end up before a cliff!)
    \begin{align*}
      \text{update}_t &= \beta\,\text{update}_{t-1}
                        + \eta\,\nabla w_{\text{look-ahead}}\\
      w_{t+1}        &= w_t - \text{update}_t
    \end{align*}
  \end{block}
\end{frame}

%==================================================================
\section{Lec 5.6: Stochastic \& Mini-batch Gradient Descent}
%==================================================================

\begin{frame}[shrink=10]{Stochastic Gradient Descent (SGD)}
  \begin{block}{Consequence of Full-batch GD}
    Full-batch requires \textbf{millions of calculations} per update step
    (one pass over all data before any parameter update).
  \end{block}

  \medskip

  \begin{algorithm}[H]
  \begin{algorithmic}[1]
    \REQUIRE $w,\,b,\,\eta$,\;\texttt{max\_epoch}
    \FOR{$i$ in \texttt{range(max\_epoch)}}
      \STATE $dw,\,db \leftarrow 0,\,0$
      \FOR{$x,\,y$ in \texttt{zip}$(X,Y)$}
        \STATE $dw = \texttt{grad\_w}(w,b,x,y)$;\quad
               $db = \texttt{grad\_b}(w,b,x,y)$
        \STATE $w \leftarrow w-\eta\,dw$;\quad $b \leftarrow b-\eta\,db$
      \ENDFOR
    \ENDFOR
  \end{algorithmic}
  \caption{Stochastic Gradient Descent (update per point)}
  \end{algorithm}

  \medskip

  \begin{alertblock}{Stochastic $=$ Like Tossing a Coin}
    \begin{itemize}\setlength\itemsep{4pt}
      \item \textbf{No theoretical guarantee} that each step decreases
            the loss
      \item Many \textbf{oscillations} --- greedy, per-point decisions
      \item Each point pushes parameters in the direction most favourable
            to \emph{it}, ignoring all other points
      \item Data points effectively \textbf{compete} with each other
    \end{itemize}
  \end{alertblock}
\end{frame}

%------------------------------------------------------------------
\begin{frame}[shrink=10]{Mini-batch Gradient Descent (MBGD)}
  \begin{block}{Idea: Better Gradient Estimate}
    SGD estimates the gradient from \textbf{1 data point} at a time.
    MBGD averages over a \textbf{batch of $B$ points}.
  \end{block}

  \medskip

  \begin{algorithm}[H]
  \begin{algorithmic}[1]
    \REQUIRE $w,b,\eta$;\;\texttt{min\_batch\_size} $B$
    \FOR{$i$ in \texttt{range(max\_epoch)}}
      \STATE $dw,db,\texttt{num\_seen} \leftarrow 0,0,0$
      \FOR{$x,y$ in \texttt{zip}$(X,Y)$}
        \STATE $dw \mathrel{+}= \texttt{grad\_w}(w,b,x,y)$;\;
               $db \mathrel{+}= \texttt{grad\_b}(w,b,x,y)$
        \STATE $\texttt{num\_seen} \mathrel{+}= 1$
        \IF{$\texttt{num\_seen}\bmod B = 0$}
          \STATE $w\leftarrow w-\eta\,dw$;\;
                 $b\leftarrow b-\eta\,db$;\;
                 $dw,db\leftarrow 0,0$
        \ENDIF
      \ENDFOR
    \ENDFOR
  \end{algorithmic}
  \caption{Mini-batch Gradient Descent}
  \end{algorithm}

  \medskip

  \begin{columns}[T]
    \column{0.5\textwidth}
    \begin{block}{Key Terms}
      \begin{itemize}\setlength\itemsep{4pt}
        \item \textbf{Epoch}: one pass over all data
        \item \textbf{Step}: one parameter update
        \item $N$ = data points;\quad $B$ = batch size
      \end{itemize}
    \end{block}

    \column{0.46\textwidth}
    \begin{block}{Steps per Epoch}
      \centering
      \renewcommand{\arraystretch}{1.3}
      \begin{tabular}{lc}
        \toprule
        \textbf{Algorithm} & \textbf{\# steps}\\
        \midrule
        Vanilla / Batch GD & $1$\\
        Stochastic GD      & $N$\\
        Mini-Batch GD      & $N/B$\\
        \bottomrule
      \end{tabular}
    \end{block}
  \end{columns}
\end{frame}

%------------------------------------------------------------------
\begin{frame}[shrink=10]{Why MBGD Oscillates Less}
  \begin{exampleblock}{Coin-Toss Analogy}
    Tossing a coin $k$ times gives a better estimate of $P(\text{heads})$
    than a single toss.
    \[
      1\ \text{toss}:\ \tfrac{1}{2}\ \big|\ \tfrac{1}{2}
      \qquad
      k\ \text{tosses}:\ \text{know which is the better estimate}
    \]
    Higher $k$ $\Rightarrow$ better estimate $\Rightarrow$ fewer
    oscillations.\\[6pt]
    Similarly: larger mini-batch $\Rightarrow$ better gradient estimate
    $\Rightarrow$ \textbf{MBGD oscillates less than SGD}.
  \end{exampleblock}

  \medskip

  \begin{block}{Relative Oscillations (large $\to$ small)}
    \[
      \text{Stochastic GD} \;>\; \text{MBGD} \;>\; \text{Vanilla GD}
    \]
    MBGD takes relatively \emph{shorter} U-turns than pure SGD.
  \end{block}
\end{frame}

%==================================================================
\section{Lec 5.7: Tips for Learning Rate \& Momentum}
%==================================================================

\begin{frame}[shrink=10]{Choosing the Learning Rate}
  \begin{alertblock}{The Dilemma}
    Set $\eta$ high to ``blow up'' the small gradient on gentle slopes
    $\Rightarrow$ steep regions diverge!
  \end{alertblock}

  \medskip

  \begin{block}{Adaptive $\eta$ Philosophy}
    \[
      \text{Gentle slope} \Rightarrow \text{fast (large } \eta\text{)}
      \qquad
      \text{Steep slope} \Rightarrow \text{slow (small } \eta\text{)}
    \]
    \begin{itemize}\setlength\itemsep{5pt}
      \item Typical values: $\eta\in\{0.0001,\;0.001,\;0.01,\;\mathbf{0.1}\}$
      \item Fine search on \textbf{log scale}:
            $0.05\to 0.1\to 0.2\to\cdots$
      \item These are \textbf{heuristics} --- no single winning strategy
      \item High $\eta$: faster convergence, risk of overshooting\\
            Low $\eta$: slower but more precise
    \end{itemize}
  \end{block}

  \medskip

  \begin{block}{Multi-start Strategy}
    Every time we reach a minimum, save $(w_0,b_0)$.
    Take the minimum which is \emph{less lossy}.
  \end{block}
\end{frame}

%------------------------------------------------------------------
\begin{frame}[shrink=10]{Annealing the Learning Rate}
  \begin{block}{Step Decay}
    \begin{itemize}\setlength\itemsep{5pt}
      \item Halve $\eta$ after every $T$ epochs
      \item Halve $\eta$ after an epoch if validation error
            $>$ previous epoch's error
      \item Save model at epoch $t$; if epoch $t{+}1$ loss $>$
            epoch $t$ loss, restore saved model and set
            $\eta \leftarrow \eta/2$
    \end{itemize}
  \end{block}

  \medskip

  \begin{block}{Exponential Decay}
    \[
      \eta = \eta_0\,e^{-kt}
      \qquad
      (\eta_0,\;k\ \text{are hyperparameters};\;t\ \text{is step number})
    \]
  \end{block}

  \medskip

  \begin{block}{$1/t$ Decay}
    \[
      \eta = \frac{\eta_0}{1+kt}
    \]
    \textbf{Better not to use} --- introduces an extra hyperparameter
    ($1\tfrac{1}{2}$ more than needed).
  \end{block}
\end{frame}

%------------------------------------------------------------------
\begin{frame}[shrink=10]{Tips for Momentum}
  \begin{block}{Schedule (Sutskever et al.\ 2013)}
    \[
      \mu_{t+1}
        = \min\!\left(
            1 - 2^{-1-\log_2\!\left(\lfloor t/250\rfloor+1\right)},\;
            \mu_{\max}
          \right)
    \]
    \begin{itemize}\setlength\itemsep{5pt}
      \item $\mu_0 = 0.5$,\quad
            $\mu_{250} = 0.75$,\quad
            $\mu_{750} = 0.875$
      \item $\mu_{\max}$ chosen from
            $\{0.999,\;0.995,\;0.99,\;0.9,\;0\}$
    \end{itemize}
    The number of effective updates
    $\approx \mu_t + \dfrac{d\ell}{dw}$.
  \end{block}
\end{frame}

%==================================================================
\section{Lec 5.8: Line Search}
%==================================================================

\begin{frame}[shrink=10]{Line Search}
  \begin{block}{Goal}
    Find the best $\eta$ \emph{at each step} from a discrete set:
    \[
      \eta \in \{0.1,\;0.5,\;1.0,\;5.0,\;10.0\}
    \]
    \begin{itemize}\setlength\itemsep{5pt}
      \item Gentle slope $\Rightarrow$ pick larger $\eta$
      \item Steep slope  $\Rightarrow$ pick small $\eta$
    \end{itemize}
  \end{block}

  \medskip

  \begin{algorithm}[H]
  \begin{algorithmic}[1]
    \STATE best $w,b \leftarrow$ worst error
    \FOR{$\eta$ in \texttt{etas}}
      \STATE $\text{tmp}w = w - \eta\,dw$;\quad
             $\text{tmp}b = b - \eta\,db$
      \IF{$\texttt{error}(\text{tmp}w,\,\text{tmp}b) < \min\,\texttt{error}$}
        \STATE best $w \leftarrow \text{tmp}w$;\quad
               best $b \leftarrow \text{tmp}b$
      \ENDIF
    \ENDFOR
  \end{algorithmic}
  \caption{Line Search}
  \end{algorithm}

  \medskip

  \begin{alertblock}{Flip Side}
    If $k$ etas are tried: $k\times$ more computation.
    However, it avoids getting stuck in a local minimum for a particular
    learning rate --- with another rate we might have reached a better
    region.
  \end{alertblock}
\end{frame}

%------------------------------------------------------------------
\begin{frame}[shrink=10]{When Do Oscillations Happen?}
  \begin{itemize}\setlength\itemsep{8pt}
    \item Oscillations occur when the loss \emph{increases} from its
          current value
    \item Whether this happens depends on the learning rate chosen
    \item Convergence is \textbf{faster} than vanilla GD even with
          oscillations
    \item The oscillations in line search differ from those in Momentum
          and NAG --- here they arise from the discrete $\eta$ set,
          not from gradient history
  \end{itemize}
\end{frame}

%==================================================================
\section{Lec 5.9: Gradient Descent with Adaptive Learning Rate}
%==================================================================

\begin{frame}[shrink=10]{Motivation: Sparse Features}
  \begin{block}{Problem Setup}
    \[
      \nabla w^1 = (f(x)-y)\cdot f(x)\cdot(1-f(x))\cdot x^1
    \]
    If feature $x^1$ is \textbf{sparse} (mostly 0):
    $\nabla w^1 \approx 0$ for most inputs
    $\Rightarrow$ $w^1$ gets very few updates.
  \end{block}

  \medskip

  \begin{alertblock}{Issue}
    \begin{itemize}\setlength\itemsep{5pt}
      \item Sparse $\Rightarrow$ update slower;\quad
            Dense $\Rightarrow$ update faster
      \item If $x^1$ is both sparse \emph{and} important,
            we'd want to take updates to $w^1$ more seriously
      \item Very common in large NNs with thousands of features
            (e.g.\ credit fraud, education data)
    \end{itemize}
  \end{alertblock}

  \medskip

  \begin{exampleblock}{Intuition for the Fix}
    \textbf{Decay the learning rate} in proportion to update history:
    \[
      \eta \;\propto\; \frac{1}{\text{update history}}
    \]
    Frequent feature $\Rightarrow$ give $\eta\!\downarrow$;\quad
    Infrequent feature $\Rightarrow$ give $\eta\!\uparrow$.
  \end{exampleblock}
\end{frame}

%------------------------------------------------------------------
\begin{frame}[shrink=10]{AdaGrad}
  \begin{block}{AdaGrad Update Rule}
    \[
      v_t = v_{t-1} + (\nabla w_t)^2
    \]
    \[
      \boxed{w_{t+1} = w_t
        - \frac{\eta}{\sqrt{v_t+\varepsilon}}\,\nabla w_t}
    \]
    \begin{itemize}\setlength\itemsep{5pt}
      \item Frequent feature
            $\Rightarrow$ $v_t$ large
            $\Rightarrow$ effective $\eta\!\downarrow$
      \item Infrequent feature
            $\Rightarrow$ $v_t$ small
            $\Rightarrow$ effective $\eta$ stays large $\uparrow$
    \end{itemize}
  \end{block}

  \medskip

  \begin{block}{Behaviour (100 random $(x,y)$ pairs, 80\% sparse)}
    Initially all algorithms (GD, Momentum, AdaGrad) move mainly along
    the \textbf{vertical ($b$) axis}; very little movement along the
    horizontal ($w$) axis.\\[4pt]
    \textbf{Why?} Because the feature for $w$ is sparse, so $w$
    undergoes few updates; $b$ is dense and updates frequently.
  \end{block}

  \medskip

  \begin{alertblock}{Flip Side of AdaGrad}
    $\sqrt{v_t}$ grows unboundedly (accumulates \emph{all} past squared
    gradients).
    The effective $\eta$ for $b$ \textbf{decays to zero} ---
    AdaGrad gets \emph{stuck} near convergence.
  \end{alertblock}
\end{frame}

%------------------------------------------------------------------
\begin{frame}[shrink=10]{RMSProp}
  \begin{block}{Idea: Exponential Moving Average (EMA) of Squared Gradients}
    Instead of accumulating \emph{all} past squared gradients, decay
    them with $\beta > 0.9$:
    \[
      v_t = \beta\,v_{t-1} + (1-\beta)(\nabla w_t)^2
    \]
    \[
      \boxed{w_{t+1} = w_t
        - \frac{\eta}{\sqrt{v_t+\varepsilon}}\,\nabla w_t}
    \]
    \begin{itemize}\setlength\itemsep{5pt}
      \item $v_1 = (1-\beta)(\nabla w_1)^2$
      \item Still accumulates history, but \emph{not} aggressively
            (only a fraction each step)
      \item \textbf{Does not kill} the learning rate
    \end{itemize}
  \end{block}

  \medskip

  \begin{exampleblock}{Result}
    RMSProp overcomes AdaGrad's decay problem by being less aggressive.
    AdaGrad got stuck near convergence (unable to move in $b$ direction).
    RMSProp \textbf{does not kill the learning rate}.
  \end{exampleblock}
\end{frame}

%------------------------------------------------------------------
\begin{frame}[shrink=10]{Adam: Adaptive Moment Estimation}
  \begin{block}{Adam $=$ RMSProp $+$ Momentum $+$ Bias Correction}
    \begin{align*}
      m_t &= \beta_1\,m_{t-1} + (1-\beta_1)\,\nabla w_t
            &\leftarrow&\ \text{1st moment (mean / direction)}\\
      v_t &= \beta_2\,v_{t-1} + (1-\beta_2)(\nabla w_t)^2
            &\leftarrow&\ \text{2nd moment (adaptive scaling)}\\[6pt]
      \hat{m}_t &= \frac{m_t}{1-\beta_1^t},
      &\hat{v}_t &= \frac{v_t}{1-\beta_2^t}
            \quad\leftarrow\text{bias correction}\\[4pt]
      \beta_1 &= 0.9,\quad \beta_2 = 0.99
    \end{align*}
    \[
      \boxed{w_{t+1} = w_t
        - \frac{\eta}{\sqrt{\hat{v}_t+\varepsilon}}\cdot\hat{m}_t}
    \]
  \end{block}
\end{frame}

%------------------------------------------------------------------
\begin{frame}[shrink=10]{Why Bias Correction in Adam?}
  \begin{block}{Derivation via EMA}
    Treat stochastic gradients $g_t = \nabla w_t$ as random variables.
    We want the update to be close to $\mathbb{E}[g_t]$.
    \[
      m_t = (1-\beta)\sum_{i=1}^{t}\beta^{t-i}\,g_i
    \]
    \[
      \mathbb{E}[m_t]
        = \mathbb{E}[g_t]\cdot(1-\beta)\sum_{i=1}^{t}\beta^{t-i}
        = \mathbb{E}[g_t]\cdot(1-\beta^t)
    \]
    (Assumes all $g_i$ from the same distribution,
    i.e.\ $\mathbb{E}[g_i]=\mathbb{E}[g_t]$.)

    \medskip

    We \textbf{want} $\mathbb{E}[m_t] = \mathbb{E}[g_t]$,
    so divide by $(1-\beta^t)$:
    \[
      \hat{m}_t = \frac{m_t}{1-\beta_1^t}
      \quad\Rightarrow\quad
      \mathbb{E}[\hat{m}_t] = \mathbb{E}[g_t]
    \]
    The same argument applies to $\hat{v}_t = v_t/(1-\beta_2^t)$.
  \end{block}
\end{frame}

%------------------------------------------------------------------
\begin{frame}[shrink=10]{AdamW (Transformer Variant)}
  \begin{block}{AdamW: Weight Decay via the Gradient System}
    In standard Adam: $g_t = \nabla_w\mathcal{L}(w_t)$.\\[4pt]
    In \textbf{AdamW} (used in Transformers), weight decay is added
    directly to the gradient:
    \[
      g_t \leftarrow \nabla_w\mathcal{L}(w_t) + \lambda\,w_t
    \]
    This yields:
    \[
      w_{t+1}
        = w_t
        - \frac{\eta\,\hat{m}_t}{\sqrt{\hat{v}_t+\varepsilon}}
        - \eta\lambda\,w_t
    \]
    The extra $-\eta\lambda w_t$ term gives true
    \textbf{weight-decay} behaviour: it optimises the loss
    and independently shrinks the weights.
  \end{block}

  \medskip

  \begin{exampleblock}{Adam vs AdamW}
    In Adam, $\lambda w_t$ goes inside the moment estimation,
    changing the adaptive scaling.
    In AdamW, $\lambda w_t$ is applied \emph{after} the
    adaptive step --- cleaner decoupling.
  \end{exampleblock}
\end{frame}

%------------------------------------------------------------------
\begin{frame}[shrink=10]{Optimizer Comparison}
  \begin{block}{Hierarchy of Preference (in practice)}
    \centering
    \renewcommand{\arraystretch}{1.35}
    \begin{tabular}{ll}
      \toprule
      \textbf{Optimizer} & \textbf{Notes}\\
      \midrule
      GD                 & Baseline\\
      GD + Momentum      & Better for gentle slopes\\
      NAG                & ``Look before you leap''\\
      AdaGrad            & Adaptive $\eta$; kills learning rate over time\\
      RMSProp            & Fixes AdaGrad decay via EMA\\
      \textbf{Adam}      & \textbf{Default choice}
                           $(\beta_1{=}0.9,\,\beta_2{=}0.999,\,
                           \varepsilon{=}10^{-8})$\\
      \bottomrule
    \end{tabular}
  \end{block}

  \medskip

  \begin{alertblock}{Caveats}
    \begin{itemize}\setlength\itemsep{5pt}
      \item Adam does \emph{not} actually converge in some real-world
            cases --- some recent work confirms this
      \item For sequence generation: $\eta = 0.001$ or $0.0001$
            work best
      \item SGD + Momentum (Nesterov/Classical) with simple annealing
            also works well in practice
      \item Adam \emph{might} be the best choice, but is not
            guaranteed to be!
    \end{itemize}
  \end{alertblock}
\end{frame}

%==================================================================
\section{Lec 6.2: Linear Algebra --- Basis \& Orthonormal Vectors}
%==================================================================

\begin{frame}[shrink=10]{Basis \& Linear Independence}
  \begin{block}{Basis}
    A set of vectors $\in\mathbb{R}^n$ is a \textbf{basis} if they are
    linearly independent and every vector $\in\mathbb{R}^n$ can be
    expressed as a linear combination of them.
  \end{block}

  \medskip

  \begin{block}{Linearly Independent Vectors}
    $v_1,\ldots,v_n$ are linearly independent iff the only solution to
    \[
      c_1 v_1 + c_2 v_2 + \cdots + c_n v_n = \mathbf{0}
    \]
    is $c_1 = c_2 = \cdots = c_n = 0$.
  \end{block}

  \medskip

  \begin{exampleblock}{Example in $\mathbb{R}^2$}
    $x=[1,0]^T$ and $y=[0,1]^T$ are the standard unit vectors.
    Any $[a,b]^T \in \mathbb{R}^2$ equals $a[1,0]^T + b[0,1]^T$.\\[4pt]
    \textbf{Nothing is sacrosanct} about this choice --- we could
    equally use $[2,3]^T$ and $[5,7]^T$ as basis vectors.
    In general, any $n$ linearly independent vectors in $\mathbb{R}^n$
    form a basis.
  \end{exampleblock}
\end{frame}

%------------------------------------------------------------------
\begin{frame}[shrink=10]{Orthonormal Basis}
  \begin{block}{Representing Any Vector}
    \[
      z = \alpha_1 u_1 + \alpha_2 u_2 + \cdots + \alpha_n u_n
    \]
    Solving for the $\alpha_i$ via Gaussian elimination:
    $O(n^3)$.
  \end{block}

  \medskip

  \begin{block}{Orthonormal $=$ Orthogonal $+$ Normal (unit vectors)}
    \[
      u_i^T u_j = 0\ \forall\,i\neq j,
      \qquad
      u_i^T u_i = \|u_i\|^2 = 1
    \]
    With an orthonormal basis we can \textbf{directly} find each
    $\alpha_i$:
    \[
      \alpha_i = z^T u_i
    \]
    Linear complexity $O(N)$ per coefficient
    $\Rightarrow$ total $O(Nn)$ vs.\ $O(n^3)$.

    \medskip

    \textbf{Orthonormal basis is the most convenient basis!}
  \end{block}
\end{frame}

%==================================================================
\section{Lec 6.1: Eigenvalues \& Eigenvectors}
%==================================================================

\begin{frame}[shrink=10]{Eigenvalues \& Eigenvectors}
  \begin{block}{Motivation}
    When a matrix $A$ multiplies a general vector, the vector
    \textbf{strays from its direction} --- it may get rotated, scaled,
    or shortened.
    For a given square $A$ there exist \textbf{special vectors that
    refuse to stray}: they only get scaled.
  \end{block}

  \medskip

  \begin{block}{Definition}
    \[
      Ax = \lambda x
    \]
    The vector changes magnitude only, \textbf{not direction}.
    $x$ is an \textbf{eigenvector}; $\lambda$ is its
    \textbf{eigenvalue}.
  \end{block}

  \medskip

  \begin{columns}[T]
    \column{0.5\textwidth}
    \begin{block}{Theorem 1}
      Eigenvectors of $A\in\mathbb{R}^{n\times n}$ with
      \textbf{distinct eigenvalues} are
      \textbf{linearly independent} $\Rightarrow$ form a basis.
    \end{block}

    \column{0.46\textwidth}
    \begin{block}{Theorem 2}
      Eigenvectors of a \textbf{square symmetric} matrix are
      \textbf{orthogonal}.
    \end{block}
  \end{columns}

  \medskip

  \begin{exampleblock}{Applications}
    PCA, Autoencoders, SVD word vectors, PageRank (limelight),
    spectral graph theory.
  \end{exampleblock}
\end{frame}

%------------------------------------------------------------------
\begin{frame}[shrink=10]{Why Use Eigenvectors as a Basis?}
  \begin{block}{Key Motivation for PCA}
    A square symmetric matrix gives a very convenient basis:
    its eigenvectors are \textbf{orthogonal} (Thm 2) and form a
    \textbf{linearly independent set} (Thm 1).
  \end{block}

  \medskip

  \begin{block}{Optimisation Connection}
    For square symmetric $n\times n$ matrix $A$:
    \begin{align*}
      \max_x\ x^TAx \quad &\text{s.t.}\ \|x\|=1
      &\Rightarrow&\ \text{eigenvector of \textbf{largest} eigenvalue}\\
      \min_x\ x^TAx \quad &\text{s.t.}\ \|x\|=1
      &\Rightarrow&\ \text{eigenvector of \textbf{smallest} eigenvalue}
    \end{align*}
    \textit{Proof:} Lagrangian $L = x^TAx - \lambda(x^Tx-1)$;
    $\;\tfrac{dL}{dx} = 2Ax - 2\lambda x = 0
    \Rightarrow Ax = \lambda x$.
  \end{block}
\end{frame}

%------------------------------------------------------------------
\begin{frame}[shrink=10]{Dominant Eigenvalue \& Steady State}
  \begin{block}{Stochastic Matrix}
    A matrix $M$ is \textbf{stochastic} if all entries are positive
    and every \emph{column} sums to 1.
  \end{block}

  \medskip

  \begin{theorem}[Dominant Eigenvalue of a Stochastic Matrix]
    The largest (dominant) eigenvalue of a stochastic matrix is
    $\mathbf{1}$.
  \end{theorem}

  \medskip

  \begin{theorem}
    If $A$ is $n\times n$ with dominant eigenvalue $d_1$, then
    $Ax_0,\,A^2x_0,\,A^3x_0,\ldots$ approaches a
    \textbf{multiple of the dominant eigenvector} of $A$.
    In general: $A^p x_0 = k|d_1|^p\,e_{d_1}$.
  \end{theorem}

  \medskip

  \begin{block}{Sequence Behaviour}
    \begin{itemize}\setlength\itemsep{5pt}
      \item $|\lambda|>1$ $\Rightarrow$ sequence \textbf{explodes}
      \item $|\lambda|<1$ $\Rightarrow$ sequence \textbf{vanishes}
      \item $|\lambda|=1$ $\Rightarrow$ sequence reaches a
            \textbf{steady state}
    \end{itemize}
  \end{block}
\end{frame}

%------------------------------------------------------------------
\begin{frame}[shrink=10]{Stochastic Matrix Example: Two-Restaurant Problem}
  \begin{block}{Setup}
    \[
      V_{(0)} = \begin{bmatrix}k_1\\k_2\end{bmatrix},
      \qquad
      M = \begin{bmatrix}p & 1-q\\ 1-p & q\end{bmatrix}
    \]
    Fraction $p$ of Chinese-restaurant students stay the next day;
    $(1-p)$ switch.
    Fraction $q$ of Mexican-restaurant students stay;
    $(1-q)$ switch.\\[4pt]
    $V_{(n)} = M^n V_{(0)}$ ---
    will the population reach a \textbf{steady state}?
  \end{block}

  \medskip

  \begin{block}{Answer via Eigenvalues}
    $M$ is stochastic $\Rightarrow$ dominant eigenvalue $= 1$
    $\Rightarrow$ $V_{(n)} \to k\cdot e_d$ (a multiple of the dominant
    eigenvector).
    \textbf{Population becomes constant after some step $n$.}
  \end{block}
\end{frame}

%==================================================================
\section{Lec 6.3: Eigenvalue Decomposition (EVD)}
%==================================================================

\begin{frame}[shrink=10]{Eigenvalue Decomposition}
  \begin{block}{Construction}
    Let $u_1,\ldots,u_n$ be eigenvectors of $A$ and
    $d_1,\ldots,d_n$ the corresponding eigenvalues.
    Let $U$ have columns $u_1,\ldots,u_n$:
    \[
      AU = [Au_1\ \cdots\ Au_n]
         = [d_1u_1\ \cdots\ d_nu_n]
         = U\Lambda
    \]
    \[
      \boxed{A = U\Lambda U^{-1}}
      \qquad\text{(Eigenvalue Decomposition)}
    \]
    \[
      \Lambda = U^{-1}AU
      \qquad\text{(Diagonalisation of }A\text{)}
    \]
  \end{block}

  \medskip

  \begin{block}{Symmetric Case}
    If $A$ is symmetric: eigenvectors are orthogonal and normalised,
    so $U^TU = I \Rightarrow U^{-1} = U^T$:
    \[
      A = U\Lambda U^T
    \]
    Sequence $x_0,Ax_0,A^2x_0,\ldots$ in terms of EVD:
    $U\Lambda^n U^T x_0$ ---
    explodes if $|\lambda|>1$, steady if $|\lambda|=1$,
    vanishes if $|\lambda|<1$.
  \end{block}
\end{frame}

%==================================================================
\section{Lec 6.4: PCA \& Its Interpretations}
%==================================================================

\begin{frame}[shrink=10]{PCA: Motivation}
  \begin{block}{Goal}
    Represent data using \textbf{fewer dimensions} such that:
    \begin{itemize}\setlength\itemsep{5pt}
      \item Data has \textbf{high variance} along these dimensions
      \item Dimensions are \textbf{linearly independent} (uncorrelated)
      \item Even better: dimensions are \textbf{orthogonal}
            (most convenient basis)
    \end{itemize}
  \end{block}

  \medskip

  \begin{block}{Covariance}
    \[
      c_{ij}
        = \frac{1}{m}\sum_{l=1}^{m}(x_{li}-\mu_i)(x_{lj}-\mu_j)
        = \frac{1}{m}(X^TX)_{ij}
        \quad(\mu_i=0\ \text{for zero-mean data})
    \]
    \begin{itemize}\setlength\itemsep{4pt}
      \item $+$cov: both increase/decrease together
      \item $-$cov: one increases as the other decreases
      \item $\approx 0$: no linear relationship
    \end{itemize}
  \end{block}

  \begin{exampleblock}{Correlation}
    Normalised covariance; direction $-1$ to $+1$ and strength.
    $+1$: perfect positive; $0$: no linear relationship;
    $-1$: perfect negative.
    High correlation $\Rightarrow$ one variable is redundant.
  \end{exampleblock}
\end{frame}

%------------------------------------------------------------------
\begin{frame}[shrink=10]{PCA: Mathematical Formulation}
  \begin{block}{Setup}
    Data matrix $X\in\mathbb{R}^{m\times n}$ (zero-mean, unit variance).
    Let $P=[p_1\;\cdots\;p_n]$ be an $n\times n$ orthonormal matrix.
    \[
      \hat{X} = XP
      \qquad\text{(projection onto new basis)}
    \]
    For an orthonormal basis: $\alpha_{ij} = x_i^T p_j$.
  \end{block}

  \medskip

  \begin{block}{Covariance of Transformed Data}
    \textbf{Theorem}: If $X$ has zero-mean columns then so does
    $\hat{X}=XP$.\\[4pt]
    \textbf{Theorem}: $X^TX$ is symmetric.\\[6pt]
    \[
      \frac{1}{m}\hat{X}^T\hat{X}
        = \frac{1}{m}P^T X^T X P
        = P^T\Sigma P,
      \qquad
      \Sigma = \tfrac{1}{m}X^TX
    \]
    We \textbf{want} $P^T\Sigma P = D$ (diagonal):
    \[
      \bigl(\tfrac{1}{m}\hat{X}^T\hat{X}\bigr)_{ij}
        = 0\ (i\neq j)
        \text{ (zero covariance)},\quad
      \neq 0\ (i=j)
      \text{ (non-zero variance)}
    \]
  \end{block}
\end{frame}

%------------------------------------------------------------------
\begin{frame}[shrink=10]{PCA: The Answer}
  \begin{block}{Which $P$ Diagonalises $\Sigma$?}
    The matrix $P$ whose columns are the
    \textbf{eigenvectors of $\Sigma = \frac{1}{m}X^TX$}:
    \[
      \Sigma = P\Lambda P^T
      \;\Rightarrow\;
      P^T\Sigma P = \Lambda\text{ (diagonal)}
    \]
    The new basis $P$ is the \textbf{eigenvector basis of $X^TX$}.
  \end{block}

  \medskip

  \begin{block}{Why Is This a Good Basis?}
    \begin{itemize}\setlength\itemsep{5pt}
      \item Eigenvectors of $X^TX$ are \textbf{linearly independent}
            (Thm 1)
      \item Eigenvectors of $X^TX$ are \textbf{orthogonal}
            (since $X^TX$ is symmetric; Thm 2)
      \item This is \textbf{Principal Component Analysis (PCA)}:
            transforms data to a basis with low covariance and high
            variance dimensions
      \item In practice, keep only the \textbf{top-$k$} dimensions
            (high eigenvalue $\Rightarrow$ high variance)
    \end{itemize}
  \end{block}
\end{frame}

%==================================================================
\section{Lec 6.5: PCA --- Interpretation 2 (Reconstruction Error)}
%==================================================================

\begin{frame}[shrink=10]{PCA Interpretation 2: Minimising Reconstruction Error}
  \begin{block}{Setup}
    Given $n$ orthonormal vectors $p_1,\ldots,p_n$ we can represent
    each $x_i$ exactly:
    \[
      x_i = \sum_{j=1}^{n}\alpha_{ij}\,p_j
    \]
    Keep only the \textbf{top-$k$ dimensions}:
    \[
      \hat{x}_i = \sum_{j=1}^{k}\alpha_{ij}\,p_j
    \]
    Select $p_j$'s to minimise reconstruction error:
    \[
      e = \sum_{i=1}^{m}(x_i-\hat{x}_i)^T(x_i-\hat{x}_i)
    \]
  \end{block}

  \medskip

  \begin{block}{Result (using orthonormality)}
    After simplification:
    \[
      e = \sum_{j=k+1}^{n} p_j^T\Sigma\,p_j
    \]
    Minimising $e$ $\Rightarrow$ select $p_{k+1},\ldots,p_n$ as the
    \textbf{smallest eigenvectors} of $\Sigma$, i.e.\
    \textbf{keep the top-$k$ eigenvectors}.
    The discarded $n-k$ directions are those with minimum variance.
  \end{block}
\end{frame}

%------------------------------------------------------------------
\begin{frame}[shrink=10]{PCA Interpretation 2: Worked Example}
  \begin{exampleblock}{Example in $\mathbb{R}^2$ (data along $y=x$)}
    New basis:
    $u_1 = \bigl[\tfrac{1}{\sqrt{2}},\tfrac{1}{\sqrt{2}}\bigr]^T$
    and
    $u_2 = \bigl[-\tfrac{1}{\sqrt{2}},\tfrac{1}{\sqrt{2}}\bigr]^T$.\\[6pt]
    Consider $x = [3.3,\ 3]$:
    \[
      \alpha_1 = x^Tu_1 = \frac{6.3}{\sqrt{2}},
      \qquad
      \alpha_2 = x^Tu_2 = \frac{0.3}{\sqrt{2}}
    \]
    \textbf{Perfect reconstruction (2-D):}
    $\hat{x} = \alpha_1 u_1 + \alpha_2 u_2 = [3.3,\ 3]$.\\[6pt]
    \textbf{Using top-1 only} (ignoring low-variance $u_2$):
    \[
      \hat{x} = \alpha_1 u_1 = [3.15,\ 3.15]
      \quad\text{(minimum reconstruction error)}
    \]
    Small $\alpha_2$ $\Rightarrow$ discarding $u_2$ introduces little
    error.
  \end{exampleblock}

  \medskip

  \begin{block}{Key Idea}
    PCA uses only the top-$k$ eigenvectors:
    \begin{enumerate}\setlength\itemsep{4pt}
      \item Covariance between new dimensions is 0
            (diagonalise covariance matrix)
      \item Discarded dimensions are those with minimum variance
            (minimum reconstruction error)
    \end{enumerate}
  \end{block}
\end{frame}

%==================================================================
\section{Lec 6.6: PCA --- Interpretation 3 (Variance)}
%==================================================================

\begin{frame}[shrink=10]{PCA Interpretation 3: Maximising Variance}
  \begin{block}{Variance along the $i$-th Dimension}
    The $i$-th dimension of the transformed data is $\hat{X}_i = Xp_i$.
    Its variance:
    \[
      \frac{\hat{X}_i^T\hat{X}_i}{m}
        = \frac{1}{m}p_i^T X^T X p_i
        = \frac{1}{m}p_i^T\lambda_i p_i
        = \frac{\lambda_i}{m}
    \]
    (since $p_i$ is an eigenvector of $X^TX$ with eigenvalue
    $\lambda_i$, and $p_i^Tp_i = 1$).
  \end{block}

  \medskip

  \begin{block}{Conclusion}
    Variance along dimension $i$ is \textbf{scaled by eigenvalue
    $\lambda_i$}.\\[4pt]
    Discarding low-eigenvalue dimensions
    $\Rightarrow$ discarding \textbf{lowest-variance} (most noisy,
    least informative) dimensions.
    $\Rightarrow$ We did the right thing!
  \end{block}

  \medskip

  \begin{block}{Quick Summary: Three Interpretations of PCA}
    \begin{enumerate}\setlength\itemsep{4pt}
      \item \textbf{Covariance}: covariance between new dimensions is
            \textbf{minimised} (diagonalise $\Sigma$)
      \item \textbf{Variance}: dimensions are chosen to have
            \textbf{high variance}
      \item \textbf{Compression}: data representable with
            \textbf{fewer dimensions}
            (minimum reconstruction error)
    \end{enumerate}
  \end{block}
\end{frame}

%------------------------------------------------------------------
\begin{frame}[shrink=10]{PCA in Practice: Eigenfaces}
  \begin{block}{Setup}
    $X\in\mathbb{R}^{m\times 10K}$
    ($100\times100$ images $= 10\,\text{K}$ dimensions).\\[4pt]
    Retain \textbf{top 50--100 eigenvectors} of $X^TX$:
    \begin{itemize}\setlength\itemsep{5pt}
      \item $X^TX$ is $n\times n$ ($n=10\,\text{K}$);
            each eigenvector is $n$-dimensional
      \item Convert each eigenvector to a $100\times100$ image
            $\Rightarrow$ \textbf{eigenface}
      \item These form a basis for any face in the dataset
    \end{itemize}
  \end{block}

  \medskip

  \begin{block}{Benefits}
    \begin{itemize}\setlength\itemsep{5pt}
      \item Any face $=$ linear combination of eigenfaces
      \item Store $p_1,\ldots,p_k$ once; then only scalars
            $\alpha_{i1},\ldots,\alpha_{ik}$ per image
      \item \textbf{Data compression}: significant storage savings with
            little quality loss
      \item Multiple images of the same person (different lighting,
            contrast) are \emph{closer} in the new space
    \end{itemize}
  \end{block}
\end{frame}

%==================================================================
%  PCA SUMMARY SLIDE (standalone section)
%==================================================================
\section{PCA: Complete Summary}

\begin{frame}[shrink=10]{PCA: Complete Summary --- Setup \& Goal}
  \begin{block}{The Problem We Solve}
    Given $m$ data points in $\mathbb{R}^n$ (zero-mean),
    find a new \textbf{orthonormal basis} of size $k\ll n$ such that:
    \begin{enumerate}\setlength\itemsep{5pt}
      \item Dimensions are \textbf{uncorrelated} (zero off-diagonal
            covariance)
      \item Data exhibits \textbf{maximum variance} in these directions
      \item \textbf{Reconstruction error} from the $k$-dim
            approximation is minimised
    \end{enumerate}
  \end{block}

  \medskip

  \begin{block}{One Algorithm, Three Views}
    \centering
    \renewcommand{\arraystretch}{1.3}
    \begin{tabular}{lcl}
      \toprule
      \textbf{Interpretation} & \textbf{Optimises} & \textbf{By}\\
      \midrule
      Decorrelation  & min off-diagonal covariance &
                       diagonalise $\Sigma$\\
      Max Variance   & max variance per dimension  &
                       keep top eigenvalues\\
      Min Error      & min reconstruction error    &
                       discard small eigenvectors\\
      \bottomrule
    \end{tabular}\\[6pt]
    \textbf{All three views give the same answer:}
    use the eigenvectors of $\Sigma = \tfrac{1}{m}X^TX$.
  \end{block}
\end{frame}

%------------------------------------------------------------------
\begin{frame}[shrink=10]{PCA: Complete Summary --- The Algorithm}
  \begin{block}{Step-by-Step}
    \begin{enumerate}\setlength\itemsep{7pt}
      \item \textbf{Centre} the data: $x_i \leftarrow x_i - \mu$
      \item \textbf{Compute} covariance matrix
            $\Sigma = \tfrac{1}{m}X^TX\in\mathbb{R}^{n\times n}$
      \item \textbf{Eigendecompose}: $\Sigma = P\Lambda P^T$
            (columns of $P$ are eigenvectors; diagonal of $\Lambda$
            are eigenvalues $\lambda_1\geq\lambda_2\geq\cdots\geq\lambda_n$)
      \item \textbf{Select} top-$k$ eigenvectors
            $P_k = [p_1,\ldots,p_k]\in\mathbb{R}^{n\times k}$
      \item \textbf{Project}: $\hat{X} = XP_k\in\mathbb{R}^{m\times k}$
    \end{enumerate}
  \end{block}

  \medskip

  \begin{exampleblock}{Key Properties of the Transformed Data}
    \begin{itemize}\setlength\itemsep{5pt}
      \item $\hat{X}$ has \textbf{zero mean} (inherited from $X$)
      \item Covariance of $\hat{X}$:
            $\tfrac{1}{m}\hat{X}^T\hat{X} = \Lambda_k$
            (diagonal $\Rightarrow$ uncorrelated dimensions)
      \item Variance along $i$-th new dimension
            $= \lambda_i / m$
      \item Fraction of variance retained
            $= \sum_{i=1}^{k}\lambda_i / \sum_{i=1}^{n}\lambda_i$
    \end{itemize}
  \end{exampleblock}
\end{frame}

%------------------------------------------------------------------
\begin{frame}[shrink=10]{PCA: Complete Summary --- Why Eigenvectors?}
  \begin{columns}[T]
    \column{0.48\textwidth}
    \begin{block}{Theorem 1}
      Eigenvectors of $A$ with \textbf{distinct eigenvalues} are
      \textbf{linearly independent}.\\[4pt]
      $\Rightarrow$ They form a basis.
    \end{block}

    \column{0.48\textwidth}
    \begin{block}{Theorem 2}
      Eigenvectors of a \textbf{symmetric} matrix are
      \textbf{orthogonal}.\\[4pt]
      $\Rightarrow$ $\Sigma = \tfrac{1}{m}X^TX$ is symmetric,
      so its eigenvectors are orthogonal.
    \end{block}
  \end{columns}

  \bigskip

  \begin{block}{Optimisation View}
    \[
      \max_{p:\|p\|=1}\ p^T\Sigma\,p \;=\; \lambda_{\max}
      \quad\text{(achieved by top eigenvector)}
    \]
    Each subsequent direction maximises remaining variance subject to
    being orthogonal to all previous directions.
  \end{block}

  \bigskip

  \begin{alertblock}{Discard Low-Eigenvalue Directions Because\ldots}
    \begin{itemize}\setlength\itemsep{4pt}
      \item Low eigenvalue $\Rightarrow$ low variance $\Rightarrow$
            mostly noise
      \item Low eigenvalue $\Rightarrow$ small $\alpha_{ij}$
            $\Rightarrow$ small reconstruction error if dropped
    \end{itemize}
  \end{alertblock}
\end{frame}

%------------------------------------------------------------------
\begin{frame}[shrink=10]{PCA: Complete Summary --- Connections \& Applications}
  \begin{block}{Connections to Other Concepts}
    \begin{itemize}\setlength\itemsep{6pt}
      \item \textbf{EVD}: $\Sigma = P\Lambda P^T$ is exactly EVD of a
            symmetric matrix; $P^{-1} = P^T$
      \item \textbf{Orthonormal basis}: eigenvectors of a symmetric
            matrix are orthogonal $\Rightarrow$ finding coordinates is
            $O(N)$ per dimension
      \item \textbf{SVD}: PCA of $X$ is equivalent to SVD of $X$;
            eigenvectors of $X^TX$ are the right singular vectors
      \item \textbf{Stochastic matrix}: dominant eigenvector gives
            the steady state distribution (same maths, different context)
    \end{itemize}
  \end{block}

  \medskip

  \begin{block}{Applications}
    \begin{itemize}\setlength\itemsep{5pt}
      \item \textbf{Eigenfaces}: face recognition using top-100
            eigenvectors of $X^TX$; any face is a linear combination
            of eigenfaces
      \item \textbf{Data compression}: store $k$ scalars per image
            instead of $n=10\,\text{K}$ pixels
      \item \textbf{Pre-processing for classifiers}: uncorrelated,
            high-variance features improve generalisation
      \item \textbf{Noise reduction}: discarding low-variance dimensions
            removes uninformative noise
    \end{itemize}
  \end{block}
\end{frame}

%==================================================================
\section{Summary}
%==================================================================

\begin{frame}[shrink=10]{Summary}
  \begin{columns}[T]
    \column{0.48\textwidth}
    \begin{block}{Optimization Algorithms}
      \begin{itemize}\setlength\itemsep{4pt}
        \item \textbf{Momentum}: exponentially weighted gradient history;
              fast on gentle slopes; U-turns but faster than vanilla GD
        \item \textbf{NAG}: ``look before you leap''; checks future slope
        \item \textbf{SGD / MBGD}: stochastic; fewer calculations;
              mini-batch reduces oscillations
        \item \textbf{AdaGrad}: adaptive $\eta$; fixes sparse features;
              kills $\eta$ via unbounded accumulation
        \item \textbf{RMSProp}: EMA decay; fixes AdaGrad
        \item \textbf{Adam}: RMSProp + Momentum + bias correction;
              \textbf{default choice}
        \item \textbf{AdamW}: Adam + weight decay in gradient
              (Transformers)
      \end{itemize}
    \end{block}

    \column{0.48\textwidth}
    \begin{block}{Linear Algebra for DL}
      \begin{itemize}\setlength\itemsep{4pt}
        \item \textbf{Basis}: linearly independent spanning set
        \item \textbf{Orthonormal basis}: most convenient; $O(N)$
              projection
        \item \textbf{Eigenvectors}: $Ax=\lambda x$; symmetric
              $\Rightarrow$ orthogonal
        \item \textbf{EVD}: $A=U\Lambda U^{-1}$; symmetric:
              $A=U\Lambda U^T$
        \item \textbf{Stochastic matrix}: dominant eigenvalue is 1;
              steady state exists
        \item \textbf{PCA}: 3 interpretations: decorrelation,
              max variance, min reconstruction error
        \item \textbf{Eigenfaces}: PCA for face compression
      \end{itemize}
    \end{block}
  \end{columns}

  \medskip

  \begin{center}
    \large\textit{%
      ``Gentle slope $\to$ slow GD, fast momentum;
       Steep slope $\to$ fast GD, large gradient, careful $\eta$.''}
  \end{center}
\end{frame}

%------------------------------------------------------------------
\begin{frame}
  \Huge{\centerline{Thank You}}
  \vspace{0.6cm}
  \large{\centerline{Questions?}}
\end{frame}

\end{document}